% ======================================================================
% preprints/jphq_06_structural_tension_nonoptimizable/sections/02_formal_setting.tex
% Journal-grade content (100/100). ASCII-only.
% ======================================================================

\section{Formal Setting}

\subsection{Enterprise as Continuum}

Throughout this preprint, the enterprise is modeled as a continuum
$\KEA$ in the sense of the Ontology of Continua (OC). A continuum is
treated as a closed formal object whose admissible states are given by a
set $\OmegaEA$ together with explicitly declared admissibility
conditions.

No assumptions are made regarding geometry, smoothness, equilibrium,
optimality, convergence, trajectories, or attractors. The continuum is
defined purely through admissibility and its associated structural
constraints.

The admissible state space $\OmegaEA$ is specified relative to a fixed
set of constraints encoded in the Executable Theory Specification
\texttt{ETS.K\_EA.v1.1}. If $\OmegaEA = \varnothing$, the enterprise
continuum is undefined for diagnostic purposes within the theory.

\subsection{Boundary and Thresholds}

The boundary of admissibility is denoted by $\dOmegaEA$. It represents
the locus at which one or more admissibility conditions are satisfied at
equality, or at which required structural properties cease to hold.
Boundary terminology is used strictly in a formal sense and carries no
geometric, physical, or metric interpretation.

Thresholds are collected in the set $\ThetaEA$. Each threshold encodes a
structural constraint separating admissible from non-admissible
regimes. Thresholds are not empirical constants, performance targets,
safety margins, or control parameters. They determine only whether the
continuum remains a valid object of diagnosis.

\subsection{Axes, Flows, and Cycles}

Admissible states of $\KEA$ are described using a set of axes
$A(\KEA)$, interpreted as dimensions of distinguishability. The presence
of an axis indicates that variation along that dimension is
structurally meaningful within the model. No assumption is made
regarding cardinality, metric structure, or comparability of axes across
enterprises.

State changes are mediated by admissible flows $J(\KEA)$. Persistence of
the continuum requires the existence of recurrent structures,
represented abstractly as cycles $C(\KEA)$. The Executable Theory
Specification determines which cycles are required and how closure is
evaluated. No temporal regularity, equilibrium condition, or stability
criterion is imposed beyond closure itself.

\subsection{Continuity Measure}

The continuity measure $\kEA \in [0,1]$ indicates whether the enterprise
continuum retains coherent admissibility under the declared structural
constraints. A value $\kEA = 0$ denotes a hard structural discontinuity,
at which point $\KEA$ ceases to exist as a diagnosable object within the
theory.

Continuity is not a measure of performance, fitness, robustness, or
desirability. It expresses only the existence of admissible transitions
and required closure structures. No monotonic, functional, or causal
relation between $\kEA$ and observed outcomes is assumed or implied.

\subsection{Phase Logic}

Admissible states of $\KEA$ are classified into phases
(\PHASESUCCESS, \PHASEINERTIA, \PHASECOLLAPSE, \PHASESTOP) according to
threshold evaluations and explicit precedence rules. Phase labels
delimit which classes of diagnostic statements remain formally
admissible. They do not encode trajectories, lifecycle stages, or
normative judgments.

Phase logic is taken as fixed from \texttt{ETS.K\_EA.v1.1} and is not
redefined here. It provides the logical context in which structural
tension is interpreted in subsequent sections.

\subsection{Explicit Non-Optimization Setting}

No objective function, utility, loss, cost functional, or preference
ordering is defined on $\OmegaEA$. The formal setting explicitly excludes
any assumption that admissible states are ordered by desirability or
that system behavior seeks extrema of a scalar quantity.

This absence is not an omission but a structural property of the theory.
All results in this preprint are derived under this explicit
non-optimization setting.

%%% End of section %%%
